\section{Empirical Strategy}\label{sec:strategy}

\subsection{OLS Benchmark}

We begin with a survey-weighted OLS regression as a benchmark:
\begin{equation}\label{eq:ols}
    \text{OOP\_share}_h = \alpha + \boldsymbol{\beta} \, \mathbf{Ins}_h + \boldsymbol{\gamma} \, \mathbf{Q}_h + \boldsymbol{\delta} \, (\mathbf{Ins}_h \times \mathbf{Q}_h) + \boldsymbol{\theta} \, \mathbf{X}_h + \boldsymbol{\phi} \, \mathbf{R}_h + \varepsilon_h
\end{equation}
where $\mathbf{Ins}_h$ is a vector of insurance dummies (SIS, EsSalud; uninsured omitted), $\mathbf{Q}_h$ is a vector of consumption quintile dummies (Q2--Q5; Q1 omitted), $\mathbf{Ins}_h \times \mathbf{Q}_h$ is the full set of insurance--quintile interactions, $\mathbf{X}_h$ is a vector of household-level controls (age, sex, and education of household head, household size, children under 5, elderly 65+, rural, chronic illness), and $\mathbf{R}_h$ is a vector of 24 department fixed effects. Standard errors are clustered at the PSU level (\texttt{CONGLOME}).

We note that linear OLS on a fractional bounded outcome $\in [0,1]$ can produce fitted values outside the unit interval and the linear approximation is inappropriate at extremes \citep{Papke1996}. We report a fractional logit specification as a robustness check in Section~\ref{sec:robustness}. The OLS benchmark is retained for interpretive transparency and comparability with the quantile regression coefficients.

The insurance--quintile interaction is the primary specification, not an alternative. Insurance categories and quintile dummies are jointly correlated by design: SIS is means-tested toward the poor, and EsSalud covers formal workers in higher quintiles. Including both without interactions partially absorbs each other's variation; the interaction allows the insurance association to vary freely across the income distribution.

\subsection{Conditional Quantile Regression}

Our main specification replaces the OLS conditional mean with the conditional $\tau$-quantile:
\begin{equation}\label{eq:qr}
    Q_\tau(\text{OOP\_share}_h \mid \mathbf{Z}_h) = \alpha(\tau) + \boldsymbol{\beta}(\tau) \, \mathbf{Ins}_h + \boldsymbol{\gamma}(\tau) \, \mathbf{Q}_h + \boldsymbol{\delta}(\tau) \, (\mathbf{Ins}_h \times \mathbf{Q}_h) + \boldsymbol{\theta}(\tau) \, \mathbf{X}_h + \boldsymbol{\phi}(\tau) \, \mathbf{R}_h
\end{equation}
where $Q_\tau(\cdot)$ denotes the $\tau$-th conditional quantile function \citep{Koenker1978}. We estimate this model at $\tau = \{0.10, 0.25, 0.50, 0.75, 0.90\}$, but restrict substantive interpretation to $\tau \geq 0.50$ given the 44.1\% zero-OOP mass (Section~\ref{sec:data}). At $\tau = 0.10$ and $\tau = 0.25$, the conditional quantile is at or near zero for most covariate configurations, and coefficients are uninformative for the distributional question of interest.

The coefficient $\beta_{\text{SIS}}(\tau)$ represents the change in the conditional $\tau$-quantile of OOP share associated with SIS membership (relative to uninsured) at Q1, holding other covariates constant. The total SIS association at quintile $q$ is $\beta_{\text{SIS}}(\tau) + \delta_{\text{SIS} \times q}(\tau)$.

We emphasize two interpretive caveats. First, quantile regression coefficients represent changes in the conditional $\tau$-quantile, not average marginal effects. A household at the 10th conditional quantile is not necessarily a ``low-spending household'' in the unconditional sense; it is a household whose spending is low relative to what its covariates predict. Second, at $\tau = 0.10$ in our setting, negative coefficients for SIS may reflect that SIS beneficiaries are more likely to have zero spending (care avoidance or access to free care), not that SIS protects against high spending. The two-part model below disentangles these channels.

Survey weights (\texttt{FACTOR07}) are applied as frequency weights rescaled to preserve the original sample size while reflecting the population distribution. We acknowledge that this approach does not fully propagate the stratified multistage survey design through the quantile regression estimator. Properly design-weighted quantile regression would require a more complex implementation \citep{Koenker2005}. As a partial check, we compare weighted and unweighted estimates in Section~\ref{sec:robustness}; large discrepancies would signal that the weighting scheme materially affects conclusions.

\subsection{Two-Part Model}

Given the 44.1\% zero-OOP mass, we complement the CQR analysis with a two-part model \citep{Manning1987}:

\textbf{Part 1 (Extensive margin):} Survey-weighted probit for $\Pr(\text{OOP} > 0)$ on the same regressors as Equation~\ref{eq:ols}. We report average marginal effects (AME), not log-odds.

\textbf{Part 2 (Intensive margin):} Conditional quantile regression at $\tau = \{0.25, 0.50, 0.75, 0.90\}$ on the positive-OOP subsample ($N = 18{,}843$). This subsample is selected (households that seek care and incur costs), introducing a Heckman-type selection concern. We note this limitation but observe that the two-part model's value lies in decomposing the extensive and intensive margins, not in correcting for selection per se.

This framework addresses the zero-mass problem more directly than standard CQR. Part~1 identifies whether insurance status affects the probability of incurring any OOP spending---the participation margin. Part~2 characterizes how insurance is associated with the intensity of spending among those who do spend. The decomposition is informative regardless of selection concerns because it separates two qualitatively different margins that standard CQR conflates at lower quantiles.

We note that censored quantile regression \citep{Powell1986, Chernozhukov2002} offers a theoretically preferable alternative that treats zero as a censoring point and produces consistent estimates at all quantiles without sample selection. Implementation of the Powell (1986) or Chernozhukov--Hong (2002) three-step estimator is computationally intensive with our sample size and interaction structure, and we leave this extension for future work.

\subsection{Inference}

Standard errors are computed via cluster bootstrap at the PSU level (\texttt{CONGLOME}) with 200 replications. The cluster bootstrap resamples entire PSUs, preserving the within-cluster correlation structure.

We acknowledge that 200 bootstrap replications is below the recommended 999 for publishable inference \citep{Cameron2008}. At $\tau = 0.90$, where the effective estimation sample is the top decile of OOP share, cluster-level sample sizes are thin and bootstrap standard errors are noisier. \textbf{Significance conclusions at upper-tail quantiles ($\tau = 0.75, 0.90$) should be interpreted with caution pending replication with 999+ replications.} We report Koenker--Bassett sandwich standard errors as a comparison in Section~\ref{sec:robustness} and note that the wild cluster bootstrap \citep{Roodman2019} would provide a more robust alternative when within-cluster sample sizes at extreme quantiles are small.

\subsection{Multiple Testing}

Our analysis involves a large family of hypothesis tests: 5 quantiles $\times$ (2 insurance dummies + 8 interaction terms + controls), plus 25 robustness regressions. Without adjustment, some results will appear significant by chance under conventional $\alpha = 0.05$ thresholds. We do not apply a formal familywise correction (Romano--Wolf or Benjamini--Hochberg) in the current version, but flag this as a limitation: the selective highlighting of significant insurance--quintile interactions at specific quantiles without multiple testing adjustment risks false discovery inflation. We report joint F-tests for the full set of interactions at each quantile as a partial mitigation.

\subsection{Identification and Limitations}

Our estimates are descriptive, not causal. Three sources of endogeneity warrant explicit discussion:

\begin{enumerate}
    \item \textit{Non-random insurance assignment.} SIS enrollment is means-tested through SISFOH: eligible households self-select based on poverty status, health needs, and administrative access. EsSalud membership is tied to formal employment. Insurance status is correlated with unobserved determinants of health spending \citep{Wagstaff2010, Acharya2012}. We do not claim causal identification.

    \item \textit{Quintile endogeneity.} Consumption quintiles are constructed from \texttt{GASHOG2D}, which includes health spending. This creates a mechanical positive correlation between quintile rank and OOP share that inflates the quintile gradient throughout our results (Section~\ref{sec:data}). The interaction terms are similarly affected.

    \item \textit{SIS targeting leakage.} SIS enrollment in upper quintiles (Q4--Q5) reflects well-documented SISFOH targeting errors \citep{Cortez2008, Francke2013}. The SIS--quintile interactions at Q4--Q5 conflate insurance protection with the composition of the (mis)targeted population. We cannot decompose legitimate enrollment from leakage in our data.
\end{enumerate}

Despite these limitations, the distributional analysis provides descriptive evidence on where financial protection gaps are concentrated---information that binary CHE indicators and mean-based regressions cannot deliver.
